% Part: computability
% Chapter: computability-theory
% Section: russells-paradox

\documentclass[../../../include/open-logic-section]{subfiles}

\begin{document}

\olfileid{cmp}{thy}{rus}
\olsection{రసెల్ వైరుధ్యంతో పోలిక}

ఈ విభాగంలోని వాదనలను రసెల్ వైరుధ్యంతో పోల్చి, వాటి సారూప్యాలూ
భేదాలూ పరిశీలించడం ఉపయుక్తం:
\begin{enumerate}
\item రసెల్ వైరుధ్యం: $S=\Setabs{x}{x\notin x}$ అని ఉంచండి. అప్పుడు
  $S\in S$ అయితేనే $S\notin S$ అవుతుంది; ఇది వైరుధ్యం.

  \emph{ముగింపు:} అటువంటి సమితి~$S$ లేదు. ``సమితులన్నిటి సమితి''
  ఉందని ఊహించడం సమితి సిద్ధాంతంలోని ఇతర స్వీకృతాలతో అసంగతం.
  \sourcecorrection{OLTECOMTHY-004}{రసెల్ వైరుధ్యపు మొదటి అంశంలో ఆంగ్ల మూలం చివరి సభ్యత్వ సూత్రంలో $S$కు బదులు నిర్వచించని పెద్ద అక్షరం $X$ను రాసింది. పక్కనే ఇచ్చిన $S=\Setabs{x}{x\notin x}$ నిర్వచనం కోరినట్లుగా దాన్ని $S\notin S$గా సరిచేశాం.}

\item రసెల్ వైరుధ్యపు ఒక మార్పు: ప్రమేయాలన్నిటి సమితి నుంచి
  $\{0,1\}$కు వెళ్లే ``ప్రమేయం'' $F$ను ఇలా నిర్వచించండి:
  \[
  F(f)=
  \begin{cases}
    1 & \text{$f$ తన స్వంత నిర్వచన క్షేత్రంలో ఉండి, $f(f)=0$ అయితే} \\
    0 & \text{ఇతర సందర్భాల్లో.}
  \end{cases}
  \]
  ఇలాంటి వాదనే $F(F)=0$ అయితేనే $F(F)=1$ అవుతుందని చూపుతుంది;
  ఇది వైరుధ్యం.

  \emph{ముగింపు:} $F$ ఒక ప్రమేయం కాదు. ``ప్రమేయాలన్నిటి సమితి'' ఒక
  ప్రమేయానికి నిర్వచన క్షేత్రంగా ఉండటానికి మరీ పెద్దది.

\item వికర్ణీకరణ వాదన: $f_0$, $f_1$, \dots{} అనేది పాక్షిక గణనీయ
  ప్రమేయాల లెక్కింపు అనుకుని, $G\colon\Nat\to\{0,1\}$ను ఇలా
  నిర్వచించండి:
  \[
  G(x)=
  \begin{cases}
    1 & \text{$f_x(x)\downarrow=0$ అయితే} \\
    0 & \text{ఇతర సందర్భాల్లో.}
  \end{cases}
  \]
  $G$ గణనీయమైతే, ఏదో $k$కు అది $f_k$ అవుతుంది. కానీ అప్పుడు
  $G(k)=1$ అయితేనే $G(k)=0$ అవుతుంది; ఇది వైరుధ్యం.

  \emph{ముగింపు:} $G$ గణనీయం కాదు. సమితి సిద్ధాంతపు స్వీకృతాల
  ప్రకారం $G$ ఇప్పటికీ ఒక ప్రమేయమేనని గమనించండి; ఇక్కడ వైరుధ్యం లేదు,
  ఉన్నది ఒక స్పష్టీకరణ మాత్రమే.
\end{enumerate}

పాక్షిక ప్రమేయాలు, గణనీయ ప్రమేయాలు, పాక్షిక గణనీయ ప్రమేయాలు మొదలైన
వాటి గురించిన మాటలు గందరగోళంగా ఉండవచ్చు. $\Nat$ నుంచి $\Nat$కు వెళ్లే
పాక్షిక ప్రమేయాలన్నిటి సమితి చాలా పెద్ద వస్తు సమాహారం. వాటిలో కొన్ని
సర్వనిర్వచితాలు, కొన్ని గణనీయాలు, కొన్ని రెండూ, మరికొన్ని రెండూ కావు.
మనం ``ప్రమేయం'' అని మాత్రమే అన్నప్పుడు, ఆనవాయితీగా సర్వనిర్వచిత
ప్రమేయాన్నే ఉద్దేశిస్తామని గుర్తుంచుకోండి. అందువల్ల మనకు కిందివి ఉన్నాయి:
\begin{enumerate}
\item గణనీయ ప్రమేయాలు
\item సర్వనిర్వచితాలు కాని పాక్షిక గణనీయ ప్రమేయాలు
\item గణనీయాలు కాని ప్రమేయాలు
\item సర్వనిర్వచితాలూ గణనీయాలూ రెండూ కాని పాక్షిక ప్రమేయాలు
\end{enumerate}
దీన్ని క్రమపరచుకోవడానికి, $\Nat$ నుంచి $\Nat$కు వెళ్లే పాక్షిక
ప్రమేయాలన్నిటినీ సూచించే ఒక పెద్ద చతురస్రాన్ని గీసి, అందులో వరుసగా
సర్వనిర్వచిత ప్రమేయాలకూ గణనీయ పాక్షిక ప్రమేయాలకూ అనుగుణమైన రెండు
అతివ్యాప్త ప్రాంతాలను గుర్తించడం సహాయపడవచ్చు. ఆ చిత్రంలో ఏర్పడే ప్రతి
ప్రాంతంలోనూ ఒక వస్తువును వర్ణించగలరేమో చూడటం మంచి సాధన.

\end{document}
